% ============================================================
%  CONCLUSION GÉNÉRALE & BIBLIOGRAPHIE
% ============================================================
\chapter*{Conclusion Générale}
\addcontentsline{toc}{chapter}{Conclusion Générale}
\thispagestyle{plain}

\chapintro{Ce chapitre synthétise les réalisations du projet MecaRage et
présente les perspectives d'évolution envisagées.}

Le projet \textbf{MecaRage} représente une solution logicielle complète et
cohérente pour la modernisation de la gestion des ateliers mécaniques. À travers
quatre sprints successifs, nous avons construit une plateforme qui couvre
l'intégralité du cycle de vie d'une intervention automobile~: de la réception
initiale des symptômes jusqu'au paiement final.

\section*{Bilan des réalisations}

\begin{table}[H]
\centering
\begin{tabularx}{\textwidth}{|l|X|c|c|}
\hline
\rowcolor{gray!20}
\textbf{Sprint} & \textbf{Fonctionnalités} & \textbf{SP réalisés} & \textbf{US} \\
\hline
Sprint 1 & Authentification JWT, inscription, refresh token, profil & 10 & 4 \\[2pt]\hline
Sprint 2 & Tenants, garages, personnel, véhicules, inventaire pièces & 21 & 5 \\[2pt]\hline
Sprint 3 & Symptômes, diagnostic IA RAG, révision chef, rendez-vous & 24 & 5 \\[2pt]\hline
Sprint 4 & Atelier, facture, approbation, réparation, paiement, PDF, mobile & 62 & 13 \\[2pt]\hline
\multicolumn{2}{|r|}{\textbf{Total}} & \textbf{117 SP} & \textbf{27 US} \\
\hline
\end{tabularx}
\caption{Récapitulatif des réalisations par sprint}
\label{tab:bilan}
\end{table}

\noindent Les principaux apports techniques du projet sont~:

\begin{itemize}[noitemsep,leftmargin=1.5em]
  \item \textbf{Clean Architecture avec CQRS}~: séparation des préoccupations,
    testabilité et maintenabilité à long terme via MediatR.
  \item \textbf{Service de diagnostic IA}~: architecture RAG (BM25 + Ollama LLM local)
    produisant un pré-diagnostic structuré et contextualisé.
  \item \textbf{Intervention Lifecycle Tracker}~: suivi transverse automatique sur
    10 états, déclenché par les événements existants sans modifier le flux métier.
  \item \textbf{Gestion multi-tenant}~: isolation totale des données de chaque
    entreprise via un filtre global EF Core (\texttt{HasQueryFilter}).
  \item \textbf{Application mobile Android (Java + SQLite)}~: stockage local hors-ligne
    avec synchronisation périodique directe vers MySQL via un service FastAPI dédié.
  \item \textbf{Génération PDF (QuestPDF)}~: factures professionnelles téléchargeables
    directement depuis l'interface client web et mobile.
  \item \textbf{Infrastructure conteneurisée}~: Docker Compose supervisé par
    Prometheus et Grafana pour le monitoring en production.
\end{itemize}

\section*{Perspectives d'évolution}

\begin{enumerate}[noitemsep,leftmargin=1.5em]
  \item \textbf{Enrichissement de la base IA}~: intégration de nouvelles marques et
    modèles de véhicules dans \texttt{knowledge\_base.txt}, fine-tuning du modèle.
  \item \textbf{Alertes de réapprovisionnement}~: décrémenter automatiquement le stock
    pièces lors de la validation d'examen et alerter sous le seuil de réordre.
  \item \textbf{Calendrier interactif}~: vue mensuelle des rendez-vous via
    \texttt{@fullcalendar/angular} pour une meilleure planification.
  \item \textbf{Export comptable}~: export CSV/Excel des interventions clôturées
    sur une période pour intégration avec les outils comptables.
  \item \textbf{Notation et avis client}~: évaluation 5 étoiles + commentaire
    après la clôture d'une intervention.
  \item \textbf{SignalR}~: mise à jour en temps réel du tableau de bord web sans
    rechargement de page (actuellement polling).
  \item \textbf{Audit Log}~: enregistrement de chaque action sensible via un
    pipeline MediatR IPipelineBehavior pour la conformité et le débogage.
\end{enumerate}

\section*{Bilan personnel}

Ce projet nous a permis d'appliquer concrètement des patterns d'architecture
logicielle modernes dans un contexte métier réel et contraignant (multi-tenant,
multi-rôle, multi-plateforme). L'intégration de technologies hétérogènes ---
.NET, Angular, Python, Java, SQLite --- dans une solution cohérente a constitué
le principal défi et la principale source d'apprentissage.

\cleardoublepage

% ——————————————————————————————————————————————————————————
\chapter*{Bibliographie}
\addcontentsline{toc}{chapter}{Bibliographie}
\thispagestyle{plain}

\begin{enumerate}[label={[\arabic*]},leftmargin=1.8em]

  \item Microsoft Corporation.
    \textit{ASP.NET Core 9 Documentation}.
    Microsoft Docs, 2025.
    \url{https://docs.microsoft.com/aspnet/core}

  \item SQLite Consortium.
    \textit{SQLite — Documentation Officielle}.
    SQLite, 2025.
    \url{https://www.sqlite.org/docs.html}

  \item FastAPI Team.
    \textit{FastAPI — High Performance Python Web Framework}.
    Tiangolo, 2025.
    \url{https://fastapi.tiangolo.com}

  \item Ollama Team.
    \textit{Ollama — Run Large Language Models Locally}.
    Ollama, 2025.
    \url{https://ollama.com}

  \item Angular Team.
    \textit{Angular v21 Documentation}.
    Angular.io, 2025.
    \url{https://angular.io/docs}

  \item Fowler, M.
    \textit{Patterns of Enterprise Application Architecture}.
    Addison-Wesley Professional, 2002.
    ISBN 0-321-12521-5.

  \item Young, G.
    \textit{CQRS Documents}.
    Greg Young, 2010.
    \url{https://cqrs.files.wordpress.com/2010/11/cqrs_documents.pdf}

  \item QuestPDF Team.
    \textit{QuestPDF — Modern .NET PDF Library}.
    QuestPDF, 2025.
    \url{https://www.questpdf.com}

  \item Robertson, S. \& Robertson, J.
    \textit{Mastering the Requirements Process}. 3\textsuperscript{e} édition.
    Addison-Wesley, 2012.
    ISBN 0-321-81574-5.

  \item Pomelo Foundation.
    \textit{Pomelo.EntityFrameworkCore.MySql — GitHub}.
    2025.
    \url{https://github.com/PomeloFoundation/Pomelo.EntityFrameworkCore.MySql}

  \item Android Developers.
    \textit{Android Java — Developer Guide}.
    Google, 2025.
    \url{https://developer.android.com/guide}

\end{enumerate}
